\chapter{Work Completed}

\section{Data Collection}

The success of any data-driven project hinges on the quality and diversity of its data sources. For our tourist attraction recommendation system, robust data collection was paramount. The foundation of our recommendation system spans multiple locations, with a special focus on identifying and collecting relevant data from various cities worldwide. Our approach to data collection prioritized sources related to tourist attractions, including local tourism websites, historical archives, community feedback platforms, and mapping services such as OpenStreetMap.

We focused on gathering information about tourist attractions like historical landmarks, museums, cultural events, local eateries, and landmarks that can cater to different user preferences. Using a combination of local sources and global ones like OpenStreetMap, we aimed to create a comprehensive and diverse dataset for the recommendation system.

\subsection{Identifying Data Sources}
To build the tourist attraction recommendation system, we identified several key data sources. Primary data was collected from popular tourism websites like TripAdvisor\footnote{TripAdvisor: \url{https://www.tripadvisor.com}}, Lonely Planet\footnote{Lonely Planet: \url{https://www.lonelyplanet.com}}, and Google Reviews\footnote{Google Reviews: \url{https://www.google.com/maps}}. Additionally, OpenStreetMap\footnote{OpenStreetMap: \url{https://www.openstreetmap.org}} played a crucial role in providing geographical and missing data, allowing us to fill in gaps related to locations and other crucial details.

These platforms provide detailed information about various attractions, including user reviews, ratings, descriptions, and geographical information, which helped build an accurate and diverse dataset for tourist attractions across different cities.

\subsection{Web Scraping}
Since there was no access to pre-existing datasets specifically for some cities, we took the initiative to create our own comprehensive dataset using web scraping techniques. Selenium was used to automate the process of navigating and interacting with dynamic web pages to extract necessary data. We focused on local tourism websites, such as Bhaktapur.com\footnote{Bhaktapur.com: \url{https://www.bhaktapur.com/}}, to collect detailed information on local attractions like temples, squares, museums, and cultural events. For other cities, we leveraged OpenStreetMap to gather geographic data, including latitude, longitude, and other location-based information. This approach allowed us to build a diverse and inclusive dataset, covering multiple cities and their tourist offerings.

\section{Data Cleaning and Preprocessing}

The initial data collected from multiple sources was often inconsistent or incomplete. To ensure the system provides accurate and reliable recommendations, we underwent a thorough data cleaning and preprocessing phase. This included handling missing entries, fixing inconsistent formatting (e.g., geographic data), and removing duplicate entries.

We also utilized OpenStreetMap to search for missing values, categories, and additional information about various attractions. For data enrichment, we added relevant categories and tags based on descriptions, user reviews, and other sources to provide users with more targeted and accurate recommendations.



\section{Preliminary Model Development}
We’ve developed two recommendation models: a content-based filtering model and a collaborative filtering model. The content-based model, trained on a large dataset of 22,000 entries, performs reliably. However, the collaborative filtering model, trained on a much smaller dataset,  which limits its effectiveness and results in suboptimal performance. To improve its performance, we’ll need to gather more data and refine our preprocessing.



